% !TeX root = ../初等数学笔记.tex
% 由主文档通过 \input 载入；请编译主文档。
\setcounter{chapter}{5} % 保留原书章号。
\chapter{不等式与不等式组}
\label{chap:em-solving-inequalities}
解不等式是找出满足指定大小关系的取值；证明一个不等式则是说明它在给定范围内始终成立。两者都以顺序性质为基础，但任务不同。本章侧重解集与参数范围，普遍估计方法见第\ref{chap:em-inequality-methods}章。

\section{不等式的性质与解集运算}
若 $a<b$，则 $a+c<b+c$；乘以 $c>0$ 保持方向，乘以 $c<0$ 反向。原因是 $b-a>0$，而 $c(b-a)$ 的符号由 $c$ 决定。未知表达式作乘除因子时必须先判断符号，尤其不能在分式不等式中不加区分地乘以分母。

不等式组取各解集的交，分类得到的候选区间取并。严格不等式不包含取等端点，非严格不等式允许取等，但定义域排除的点永远不能加入。例如 $x>1$ 且 $x\le3$ 的解集为 $(1,3]$；$x<-1$ 或 $x>1$ 的解集是两个区间之并。

\section{一次与二次不等式}
一次不等式 $ax>b$ 在 $a>0$ 时为 $x>b/a$，在 $a<0$ 时为 $x<b/a$，在 $a=0$ 时根据 $0>b$ 的真假分别得到全体实数或空集。

对 $a\ne0$ 的二次式 $q(x)=ax^2+bx+c$，若 $\Delta<0$，配方
\[
 q(x)=a\left[\left(x+\frac b{2a}\right)^2-\frac{\Delta}{4a^2}\right]
\]
表明它始终与 $a$ 同号。若 $\Delta=0$，除重根外同样与 $a$ 同号，重根处为零。若有两不同根 $\alpha<\beta$，则 $q(x)=a(x-\alpha)(x-\beta)$，根外两因子同号，根内异号。这给出全部二次不等式的符号规则，而不必只靠抛物线图像记忆。

\section{多项式与分式的符号分析}
将可分解表达式写成常数与一次因子幂的乘积，并将分子零点、分母零点按大小排列。每个开区间内各因子的符号不变；越过奇重根时符号改变，越过偶重根时不变。分母零点即使同时是分子零点，也须从原定义域排除。

\begin{example}
解 $\dfrac{(x-4)^2(x+2)}{x-3}\le0$。
\end{example}
\begin{solve}
临界点为 $-2,3,4$，其中 $3$ 不允许取。除 $x=4$ 外，平方因子为正，故只需判断 $(x+2)/(x-3)$，它在 $[-2,3)$ 上非正；$x=4$ 使原式为零，也应加入。因此解集为 $[-2,3)\cup\{4\}$。偶重根可以产生孤立解，不能只记录变号点。
\end{solve}

\section{绝对值、根式与混合条件}
对 $r\ge0$，$|u|\le r$ 等价于 $-r\le u\le r$；若 $r<0$ 则无解。$|u|\ge r$ 在 $r\le0$ 时恒成立，在 $r>0$ 时等价于 $u\le-r$ 或 $u\ge r$。多个绝对值以各内部表达式的零点分段，逐段去绝对值，最终仍与该段范围相交。

根式比较须先建立定义域。求解 $\sqrt{x+2}>x$ 时，$x\ge-2$；在 $-2\le x<0$ 上右端为负，原式自动成立；在 $x\ge0$ 上可等价平方，得到 $x+2>x^2$，即 $0\le x<2$。合并为 $[-2,2)$。若直接平方并保留整个平方后的区间，将错误排除部分负数。

\section{含参数不等式与恒成立问题}
“存在 $x$ 使不等式成立”和“对所有 $x$ 都成立”具有不同量词。对 $a>0$ 的二次式，$ax^2+bx+c\ge0$ 对全部实数成立，当且仅当 $\Delta\le0$：配方后最小值为 $-\Delta/(4a)$。若 $a=0$，必须另作一次式或常数式讨论；若 $a<0$，不可能在全体实数上非负。

在闭区间上判断恒成立，则看该区间最小值。向上开口的二次式若顶点落在区间内，取顶点值；否则函数值随离顶点的距离增加，最小值在离顶点较近的端点取得。例如 $x^2-2tx+1\ge0$ 对所有 $x\in[0,1]$ 成立，$t\le0$ 时最小值为 $1$；$0<t<1$ 时为 $1-t^2>0$；$t\ge1$ 时为 $2-2t$，故总条件为 $t\le1$。

\section{例题、习题与学习检查}
\begin{example}
解 $|x-1|+|x+1|\le4$。
\end{example}
\begin{solve}
在 $x\le-1$ 时左端为 $-2x$，给出 $-2\le x\le-1$；在 $-1\le x\le1$ 时左端恒为 $2$；在 $x\ge1$ 时左端为 $2x$，给出 $1\le x\le2$。并集为 $[-2,2]$。各分段边界重复计入并不影响集合结果。
\end{solve}
\begin{enumerate}
\item 解 $(x-1)/(x+2)>0$。答：$(-\infty,-2)\cup(1,+\infty)$。
\item 求使 $x^2+2x+k\ge0$ 对所有实数 $x$ 成立的 $k$。答：$k\ge1$。
\item 解 $\sqrt{x+1}\le x-1$。答：$[3,+\infty)$；先有 $x\ge1$，再平方。
\end{enumerate}

